\documentclass[11pt,a4paper]{article}

\usepackage[margin=2.5cm]{geometry}
\usepackage{amsmath,amssymb,amsfonts}
\usepackage{algorithm}
\usepackage{algpseudocode}
\usepackage{graphicx}
\usepackage{bm}
\usepackage{cite}
\usepackage{hyperref}
\usepackage{booktabs}
\usepackage{multirow}
\usepackage{array}

\newcommand{\R}{\mathbb{R}}
\newcommand{\C}{\mathbb{C}}
\newcommand{\erfi}{\mathrm{erfi}}
\newcommand{\erf}{\mathrm{erf}}
\newcommand{\TV}{\mathrm{TV}}
\newcommand{\diag}{\mathrm{diag}}
\newcommand{\Tr}{\mathrm{Tr}}
\newcommand{\sgn}{\mathrm{sgn}}
\newcommand{\sinc}{\mathrm{sinc}}
\DeclareMathOperator*{\argmin}{arg\,min}

\title{Distributed Phased Array Forward Scatter Radar:\\
Mathematical Model and Signal Processing Algorithms\\[6pt]
\large Version 5.2 --- Oblique Motion with 3D Target Shadow Projection}

\author{D.~Huang}
\date{\today}

\begin{document}
\maketitle

\begin{abstract}
This document presents the complete mathematical model and detailed
signal processing algorithms for a distributed phased array forward
scatter radar (FSR) system.  The target is a three-dimensional rigid
body that moves along a straight line in the $x$--$y$ plane at an
arbitrary angle~$\theta$ to the baseline, with a constant
offset~$z_T$ in the $z$-direction.  Due to the oblique heading, the
target's physical extent in the body-frame $x_b$-direction maps to
different $y$-positions in the world frame; consequently, each column
of the shadow aperture experiences a \emph{distinct} effective
propagation distance $R_{\mathrm{eff},q}(t) = R(t) - X'_q\tan\theta$,
making the projected shadow silhouette time-varying and the
$z$-direction Fresnel coefficient pixel-column--dependent
($G_{p,q}(t)$ instead of $G_p(t)$).

The $K$ receiver nodes are uniformly spaced at half the Fresnel zone
radius.  A two-zone processing architecture is described.
\textbf{Zone~B} employs a phase-compensated spatial covariance matrix
to handle the time-varying DOA, with a derived formula for the
optimal number of snapshots~$M$.  \textbf{Zone~A} uses a
gradient-based iterative method with total-variation (TV) denoising
for shadow profile retrieval.  A unified noise model is used
throughout: complex AWGN on the received signal, with the FSSR
measured as the squared magnitude of the noisy total field divided by
the (perfectly known) direct-path power.
\end{abstract}

\tableofcontents
\newpage

%======================================================================
\section{System Geometry}
\label{sec:geometry}
%======================================================================

\subsection{Coordinate System}

We adopt a Cartesian system $(x,y,z)$ in which the incident plane wave
travels in the $+y$ direction.  A set of $K$ phased array receiver
nodes is positioned along the $x$-axis at $y=R_0$ (the receiver plane).
Node~$k$ is at $(\bar{X}_k,\,R_0,\,0)$, $k=0,\ldots,K{-}1$, with
$N$ elements spaced at $d=\lambda/2$ along~$x$.  Element~$n$ of
node~$k$ is at
\begin{equation}
  \mathbf{p}_{k,n} = (\bar{X}_k + n\,d,\;R_0,\;0),
  \quad n=0,\ldots,N{-}1.
  \label{eq:element_pos}
\end{equation}
All distances are normalised to $\lambda$, so $\lambda=1$ and
$d=\tfrac12$.

\subsection{Node Placement}

The $K$ nodes are uniformly spaced at half the Fresnel zone radius
evaluated at the nominal distance~$R_0$:
\begin{equation}
  \Delta X = \frac{\sqrt{R_0}}{2},
  \qquad
  \bar{X}_k = \Bigl(k - \frac{K{-}1}{2}\Bigr)\frac{\sqrt{R_0}}{2},
  \quad k=0,\ldots,K{-}1.
  \label{eq:node_positions}
\end{equation}

\subsection{Target Model --- Oblique Motion}
\label{sec:target_model}

The target is a \textbf{three-dimensional rigid body}.  Its centre
moves along a straight line in the $x$--$y$ plane at speed~$v$ and
heading angle~$\theta$ (measured from the $x$-axis toward the
$y$-axis):
\begin{equation}
  \text{Target centre at time } t:\quad
  \bigl(\,x_T(t),\;y_T(t),\;z_T\,\bigr),
  \label{eq:target_pos}
\end{equation}
where
\begin{align}
  x_T(t) &= x_0 + v_x\,t, &\quad v_x &= v\cos\theta,
  \label{eq:xT}\\
  y_T(t) &= v_y\,t,        &\quad v_y &= v\sin\theta,
  \label{eq:yT}
\end{align}
with $x_0$ the initial $x$-coordinate, $y_T(0)=0$, and $z_T$ a
constant $z$-offset.

\subsection{Time-Varying Propagation Distance}
\label{sec:R_of_t}

The nominal propagation distance from the target plane to the
receiver plane along the $y$-axis:
\begin{equation}
  R(t) = R_0 - v_y\,t = R_0 - v\sin\theta\;\!t.
  \label{eq:R_of_t}
\end{equation}

%======================================================================
\section{3D Target Shape and Shadow Projection}
\label{sec:3d_shadow}
%======================================================================

\subsection{Body-Frame Shape Function}

The target's shape is described by a binary function in the body
coordinate system $(x_b,y_b,z_b)$:
\begin{equation}
  T(x_b,y_b,z_b) =
  \begin{cases}
    1, & (x_b,y_b,z_b)\in\mathcal{V},\\
    0, & \text{otherwise},
  \end{cases}
  \label{eq:body_shape}
\end{equation}
where $\mathcal{V}$ is the target's volume.  The body frame is
oriented with:
\begin{itemize}
  \item $x_b$: along the heading direction (velocity vector
    projection in the $x$--$y$ plane),
  \item $y_b$: perpendicular to heading in the $x$--$y$ plane,
  \item $z_b$: vertical (parallel to world $z$).
\end{itemize}

\subsection{Body-Frame to World-Frame Transformation}

A point $(x_b,y_b,z_b)$ in the body frame maps to world
coordinates at time~$t$:
\begin{align}
  x_w &= x_T(t) + x_b\cos\theta - y_b\sin\theta,
  \label{eq:xw}\\
  y_w &= y_T(t) + x_b\sin\theta + y_b\cos\theta,
  \label{eq:yw}\\
  z_w &= z_T + z_b.
  \label{eq:zw}
\end{align}

\subsection{Shadow Silhouette in the Aperture Plane}
\label{sec:shadow_silhouette}

In the Fresnel diffraction framework, each body element at world
position $(x_w,y_w,z_w)$ blocks the incident plane wave (travelling
in $+y$) and produces a scattered field at the receiver.  The
propagation distance from this element to the receiver plane at
$y=R_0$ is
\begin{equation}
  R_{\rm eff}(x_b,y_b) = R_0 - y_w
  = R(t) - x_b\sin\theta - y_b\cos\theta.
  \label{eq:R_eff_general}
\end{equation}

The shadow on the $x$--$z$ plane (perpendicular to the baseline) is
the projection of the 3D body along the $y$-axis.  The \textbf{shadow
profile function} in aperture-plane coordinates $(x',z')$ relative to
the target centre is
\begin{equation}
  F(x',z') =
  \begin{cases}
    1, & \exists\;y_b:\;T(x'/\!\cos\theta + y_b\sin\theta/\!\cos\theta,
    \;y_b,\;z')=1,\\
    0, & \text{otherwise},
  \end{cases}
  \label{eq:shadow_profile}
\end{equation}
where the aperture-plane coordinates are related to the body
coordinates by
\begin{align}
  x' &= x_b\cos\theta - y_b\sin\theta, \label{eq:xprime}\\
  z' &= z_b. \label{eq:zprime}
\end{align}
For a thin target ($y_b\approx 0$):
$x'=x_b\cos\theta$, so the shadow width in~$x$ is
\textbf{foreshortened by $\cos\theta$}.

\subsection{Pixel-Dependent Effective Propagation Distance}
\label{sec:pixel_R}

\textbf{Key consequence of the 3D body:} each body element at
body-frame coordinate~$x_b$ sits at a different $y$-position in the
world frame.  For a thin target ($y_b\approx 0$), the body-frame
coordinate is $x_b = x'/\!\cos\theta$ and the world $y$-position is
$y_w = y_T(t) + x_b\sin\theta = y_T(t) + x'\tan\theta$.  Hence, the
effective propagation distance for a shadow element at aperture
coordinate~$x'$ is
\begin{equation}
  \boxed{R_{\rm eff}(x',t) = R(t) - x'\tan\theta.}
  \label{eq:R_eff_pixel}
\end{equation}

This is \textbf{both spatially varying} (different for each pixel
column~$q$) \textbf{and time-varying} (through $R(t)$).  As a result:
\begin{enumerate}
  \item The effective shadow silhouette, as weighted by the Fresnel
    propagation kernel, is \textbf{time-varying}: as $R(t)$ changes,
    the relative Fresnel-zone sizes at different pixel columns change,
    altering how each column contributes to the FSSR.
  \item The $z$-direction Fresnel coefficient $G_p$ depends on the
    pixel column~$q$ (through $R_{\rm eff,q}$), becoming
    $G_{p,q}(t)$.  The simple outer-product structure
    $\mathbf{G}\otimes\mathbf{F}$ is broken.
\end{enumerate}

For $\theta=0$: $R_{\rm eff}=R_0$ for all $x'$, recovering the
standard thin-screen model with column-independent~$G_p$.

%======================================================================
\section{FSSR Model}
\label{sec:fssr}
%======================================================================

\subsection{Continuous FSSR with Pixel-Dependent Propagation}

The Fresnel diffraction integral from a body element at aperture
coordinate $(x',z')$ to the receiver at $(\bar{X}_k,R_0,0)$ uses the
effective distance~\eqref{eq:R_eff_pixel}.  The FSSR at node~$k$:
\begin{equation}
  \boxed{
  \varepsilon_k(t) = \biggl|1
  + \iint F(x',z')\;\frac{i}{R_{\rm eff}(x',t)}\;
  \exp\!\biggl[\frac{i\pi}{R_{\rm eff}(x',t)}\Bigl(
  \bigl(\Delta x_k(t)+x'\bigr)^2
  + \bigl(z_T+z'\bigr)^2\Bigr)\biggr]
  \frac{dx'\,dz'}{\cos\theta}\biggr|^2,}
  \label{eq:fssr_continuous}
\end{equation}
where $\Delta x_k(t) = x_T(t)-\bar{X}_k$ and the factor
$1/\!\cos\theta$ accounts for the Jacobian of the coordinate
transformation from body frame to aperture plane
($dx_b = dx'/\!\cos\theta$).

\subsection{Pixel Discretisation}

The retrieval region is divided into $N_{\rm pix}=M_z\times M_x$
square pixels of side~$2\delta$ in the aperture-plane coordinates.
Pixel~$(p,q)$ has centre $(X'_q,Z'_p)$ and value
$P_{p,q}\in\{0,1\}$.  The effective propagation distance for pixel
column~$q$:
\begin{equation}
  R_{\mathrm{eff},q}(t) = R(t) - X'_q\tan\theta.
  \label{eq:R_eff_q}
\end{equation}

The FSSR becomes
\begin{equation}
  \varepsilon_k(t) \approx \biggl|1
  + \frac{1}{\cos\theta}\sum_{p=0}^{M_z-1}\sum_{q=0}^{M_x-1}
  \frac{i}{R_{\mathrm{eff},q}(t)}\;
  F_q^{(k)}(t)\;G_{p,q}(t)\;P_{p,q}\biggr|^2,
  \label{eq:fssr_pixel}
\end{equation}
where the Fresnel coefficients are now:

\subsubsection{$x$-Direction Fresnel Coefficient $F_q^{(k)}(t)$}
\begin{equation}
  F_q^{(k)}(t)
  = \sqrt{\frac{R_{\mathrm{eff},q}(t)}{4i}}\,\biggl[
    \erfi\!\biggl(\sqrt{\frac{i\pi}{R_{\mathrm{eff},q}(t)}}
    \bigl(X'_q{+}\delta{+}\Delta x_k(t)\bigr)\biggr)
  - \erfi\!\biggl(\sqrt{\frac{i\pi}{R_{\mathrm{eff},q}(t)}}
    \bigl(X'_q{-}\delta{+}\Delta x_k(t)\bigr)\biggr)\biggr].
  \label{eq:Fq}
\end{equation}

\subsubsection{$z$-Direction Fresnel Coefficient $G_{p,q}(t)$}
\begin{equation}
  G_{p,q}(t)
  = \sqrt{\frac{R_{\mathrm{eff},q}(t)}{4i}}\,\biggl[
    \erfi\!\biggl(\sqrt{\frac{i\pi}{R_{\mathrm{eff},q}(t)}}
    \bigl(Z'_p{+}\delta{+}z_T\bigr)\biggr)
  - \erfi\!\biggl(\sqrt{\frac{i\pi}{R_{\mathrm{eff},q}(t)}}
    \bigl(Z'_p{-}\delta{+}z_T\bigr)\biggr)\biggr].
  \label{eq:Gpq}
\end{equation}

\paragraph{Key change from v5.1:} $G_{p,q}$ now depends on
\textbf{both} the row index~$p$ \textbf{and} the column index~$q$
(through $R_{\mathrm{eff},q}$).  It is a full $M_z\times M_x$ matrix
rather than an $M_z$-vector.

\subsection{Combined Fresnel Coefficient Matrix}

Define the \textbf{full coefficient matrix} for node~$k$ at time~$t$:
\begin{equation}
  H_{k,p,q}(t) = \frac{i}{\cos\theta\;R_{\mathrm{eff},q}(t)}\;
  F_q^{(k)}(t)\;G_{p,q}(t)
  \;\in\;\C.
  \label{eq:H_matrix}
\end{equation}
Then
\begin{align}
  S_k(t) &= \sum_{p,q} H_{k,p,q}(t)\,P_{p,q},
  \label{eq:Sk}\\
  \varepsilon_k(t) &= |1 + S_k(t)|^2. \label{eq:eps_from_S}
\end{align}

\paragraph{Computational note.}
Since $R_{\mathrm{eff},q}$ depends only on~$q$ (not~$p$), the
$G_{p,q}$ matrix has a special structure: each column~$q$ is computed
with the \emph{same} $R$ value.  Vectorisation is still efficient:
for each column~$q$, compute $G_{0,q},\ldots,G_{M_z-1,q}$ as a
batch with $R=R_{\mathrm{eff},q}$.

%======================================================================
\section{Unified Noise Model}
\label{sec:noise}
%======================================================================

A single noise model is used in both zones.  No separate FSSR noise
is introduced.

\subsection{Complex Signal Noise}
At element~$n$ of node~$k$:
\begin{equation}
  \eta_{k,n}(t) \sim \mathcal{CN}(0,\,\sigma^2),
  \qquad
  \sigma^2 = 10^{-\mathrm{SNR}_{\mathrm{dB}}/10},
  \label{eq:noise_def}
\end{equation}
referenced to $|A_D|^2=1$.

\subsection{FSSR Measurement via Total Field}
\label{sec:fssr_measurement}

The direct-path amplitude $A_D=1$ is \textbf{perfectly known}.  After
beamforming toward broadside at node~$k$:
\begin{equation}
  \tilde{y}_k(t) = \frac{1}{N}\sum_{n=0}^{N-1} x_{k,n}(t)
  \approx 1 + S_k(t) + \tilde{\eta}_k(t),
  \label{eq:beamformed_model}
\end{equation}
where
$\tilde{\eta}_k(t) \sim \mathcal{CN}(0,\,\sigma^2/N)$
is the beamformed noise.  The FSSR is measured as
\begin{equation}
  \boxed{b_k(t) = |\tilde{y}_k(t)|^2
  = |1 + S_k(t) + \tilde{\eta}_k(t)|^2.}
  \label{eq:fssr_measured}
\end{equation}
Expanding:
\begin{equation}
  b_k(t) = \underbrace{|1{+}S_k|^2}_{\varepsilon_k}
  + \underbrace{2\,\mathrm{Re}\!\bigl[(1{+}S_k)\,
  \tilde{\eta}_k^*\bigr]}_{\text{signal-dependent}}
  + \underbrace{|\tilde{\eta}_k|^2}_{\text{quadratic}}.
  \label{eq:fssr_expanded}
\end{equation}
The effective noise on FSSR is signal-dependent and non-Gaussian.

%======================================================================
\section{Zone Classification}
\label{sec:zones}
%======================================================================

\subsection{Per-Node Fresnel Parameter}
\begin{equation}
  \xi_k(t) = \frac{|\Delta x_k(t)|}{\sqrt{R(t)}}.
  \label{eq:xi_k}
\end{equation}
\begin{equation}
  \text{Zone~B: } \xi_k(t)>\xi_{\mathrm{th}},
  \qquad
  \text{Zone~A: } \xi_k(t)\le\xi_{\mathrm{th}}.
  \label{eq:zone_def}
\end{equation}

\subsection{Sequential Zone Entry}
Node~$k$ enters Zone~A around time
$t_k^*\approx(\bar{X}_k-x_0)/v_x$.

%======================================================================
\section{Zone~B: Parameter Estimation}
\label{sec:zone_b}
%======================================================================

When $\xi_k\gg 1$, the target appears as a point source.  The
unknowns are $v_x,v_y,x_0,R_0,z_T$
(or equivalently $v,\theta,x_0,R_0,z_T$).

\subsection{Received Signal Model}
\label{sec:zone_b_signal}

At element~$n$ of node~$k$:
\begin{equation}
  x_{k,n}(t) = A_D + G_T(t)\,e^{j\psi_k(t)}\,
  e^{j\pi n\,\ell_x^{(k)}(t)} + \eta_{k,n}(t),
  \label{eq:rx_signal}
\end{equation}
where:
\begin{itemize}
  \item $A_D=1$;
  \item $G_T(t) = 2\sqrt{\pi}\,A_\Sigma / R(t)$: scattered
    amplitude;
  \item $\psi_k(t) = \pi[\Delta x_k(t)^2{+}z_T^2]/R(t)$:
    propagation phase;
  \item $\ell_x^{(k)}(t) = \Delta x_k(t)/D_k(t)$: direction cosine,
    $D_k=\sqrt{\Delta x_k^2+R(t)^2+z_T^2}$;
  \item $\eta_{k,n}\sim\mathcal{CN}(0,\sigma^2)$.
\end{itemize}

\subsection{The Time-Varying DOA Problem}
\label{sec:tv_doa}

The standard spatial covariance estimate
\begin{equation}
  \hat{\mathbf{R}}_k = \frac{1}{M}\sum_{m=1}^{M}
  \mathbf{x}_k(t_m)\,\mathbf{x}_k^H(t_m)
  \label{eq:cov_standard}
\end{equation}
assumes a stationary DOA across the $M$ snapshots.  Since
$\ell_x^{(k)}(t)$ changes with time, large~$M$ causes DOA smearing.

\subsubsection{DOA Variation Analysis}

The target's contribution to the averaged signal covariance is
\begin{equation}
  [\bar{\mathbf{R}}_{\mathrm{sig}}]_{n_1,n_2}
  = \frac{|G_T|^2}{M}\sum_{m=1}^{M}
  e^{j\pi(n_1-n_2)\ell_x(t_m)}.
  \label{eq:cov_sig}
\end{equation}
For linearly varying DOA $\ell_x(t_m)\approx\ell_0+\beta\,m\,\Delta t$:
\begin{equation}
  \frac{1}{M}\sum_{m=0}^{M-1}e^{j\pi(n_1-n_2)\beta\,m\,\Delta t}
  = \frac{\sin\!\bigl(\tfrac{\pi}{2}(n_1{-}n_2)
  \beta M\Delta t\bigr)}
  {M\sin\!\bigl(\tfrac{\pi}{2}(n_1{-}n_2)\beta\Delta t\bigr)}\;
  e^{j\phi}.
  \label{eq:cov_sum}
\end{equation}
The magnitude decays when the \emph{DOA spread}
$\Delta\ell=|\beta|\,M\,\Delta t$ becomes comparable to the array
resolution.  For the maximum baseline $(n_1{-}n_2=N{-}1)$, severe
attenuation occurs when
\begin{equation}
  M_{\max}^{\mathrm{(unc)}} \approx
  \frac{2}{(N{-}1)\,|\beta|\,\Delta t}.
  \label{eq:Mmax_uncomp}
\end{equation}

\subsubsection{Phase-Compensated Spatial Covariance}
\label{sec:phase_comp}

Given an estimate $\hat{\beta}$ of the DOA rate, define compensated
snapshots:
\begin{equation}
  \boxed{
  \tilde{x}_{k,n}(t_m) = x_{k,n}(t_m)\;
  e^{-j\pi\,n\;\!\hat{\beta}\,(t_m - t_{\mathrm{ref}})}.}
  \label{eq:phase_comp}
\end{equation}
The compensated covariance is
\begin{equation}
  \tilde{\mathbf{R}}_k
  = \frac{1}{M}\sum_{m=1}^{M}
  \tilde{\mathbf{x}}_k(t_m)\,\tilde{\mathbf{x}}_k^H(t_m).
  \label{eq:cov_compensated}
\end{equation}
If $\hat{\beta}=\beta$, the DOA is frozen at
$\ell_x(t_{\mathrm{ref}})$ for all~$m$, enabling arbitrarily
large~$M$ (up to the linearity assumption).

\subsubsection{Residual DOA Variation and the Optimal~$M$}
\label{sec:optimal_M}

After compensation, the residual DOA is
\begin{equation}
  \tilde{\ell}_x(t_m) = \ell_x(t_{\mathrm{ref}})
  + (\beta{-}\hat{\beta})(t_m{-}t_{\mathrm{ref}})
  + \gamma\,(t_m{-}t_{\mathrm{ref}})^2 + \cdots
  \label{eq:residual_doa}
\end{equation}

\paragraph{(a) Rate estimation error $\Delta\beta$:}
$|\Delta\beta|\,M\,\Delta t \le c_1/N
\;\Longrightarrow\;
M \le c_1/(N|\Delta\beta|\Delta t)$.

\paragraph{(b) Quadratic DOA variation~$\gamma$:}
With oblique motion the DOA is nonlinear
(Section~\ref{sec:doa_nonlinear}), with quadratic coefficient
\begin{equation}
  \gamma_k \approx \frac{v_x\,v_y}{R_0^2}
  + \frac{v_y^2\,(x_0{-}\bar{X}_k)}{R_0^3}.
  \label{eq:gamma_approx}
\end{equation}
Spread: $|\gamma|(M\Delta t)^2/4 \le c_2/N
\;\Longrightarrow\;
M \le 2\sqrt{c_2/N}\,/\,(|\gamma|^{1/2}\Delta t)$.

\paragraph{Combined optimal~$M$:}
\begin{equation}
  \boxed{
  M_{\mathrm{opt}} = \min\!\left(
    \frac{c_1}{N\,|\Delta\beta|\,\Delta t},\;
    \frac{2}{\sqrt{N\,|\gamma|}\;\!\Delta t},\;
    M_{\mathrm{avail}}\right).}
  \label{eq:M_opt}
\end{equation}

\paragraph{Iterative refinement:}
\begin{enumerate}
  \item \textbf{Pass~1} (no compensation): $M\le M_{\max}^{(\mathrm{unc})}$;
    estimate rough~$\hat{\beta}_k$.
  \item \textbf{Pass~2+}: compensate with~$\hat{\beta}_k$, enlarge $M$
    to~$M_{\mathrm{opt}}$, re-estimate DOA.
\end{enumerate}

\subsection{DOA Time Series and Nonlinear Model}
\label{sec:doa_nonlinear}

In the far-field approximation ($D_k\approx R(t)$):
\begin{equation}
  \ell_x^{(k)}(t) \approx \frac{\Delta x_k(t)}{R(t)}
  = \frac{x_0+v_x t-\bar{X}_k}{R_0-v_y t}.
  \label{eq:doa_ratio}
\end{equation}
Taylor expansion about $t=0$ (for $|v_y t|\ll R_0$):
\begin{align}
  \ell_x^{(k)}(t) &\approx \alpha_k + \beta_k\,t + \gamma_k\,t^2
  + \cdots, \label{eq:doa_taylor}
\end{align}
where
\begin{align}
  \alpha_k &= \frac{x_0{-}\bar{X}_k}{R_0},
  \label{eq:alpha_k}\\
  \beta_k &= \frac{v_x}{R_0}
  + \frac{v_y(x_0{-}\bar{X}_k)}{R_0^2},
  \label{eq:beta_k}\\
  \gamma_k &= \frac{v_x v_y}{R_0^2}
  + \frac{v_y^2(x_0{-}\bar{X}_k)}{R_0^3}.
  \label{eq:gamma_k}
\end{align}

\textbf{Key observation:} the DOA rate~$\beta_k$ is
\textbf{node-dependent}:
\begin{equation}
  \beta_k = \underbrace{\frac{v_x}{R_0}
  + \frac{v_y\,x_0}{R_0^2}}_a
  - \underbrace{\frac{v_y}{R_0^2}}_b\;\bar{X}_k
  \;\triangleq\; a - b\,\bar{X}_k.
  \label{eq:beta_linear_in_Xk}
\end{equation}
A linear fit of $\hat{\beta}_k$ vs.\ $\bar{X}_k$ yields slope
$\hat{b}=-v_y/R_0^2$ and intercept~$\hat{a}$.

\subsection{Doppler Analysis}
\label{sec:doppler}

The propagation phase is
$\psi_k(t)=\pi[\Delta x_k(t)^2+z_T^2]/R(t)$.  The instantaneous
Doppler:
\begin{equation}
  f_D^{(k)}(t) = \frac{1}{2\pi}\frac{d\psi_k}{dt}
  = \frac{v_x\,\Delta x_k(t)}{R(t)}
  + \frac{v_y[\Delta x_k(t)^2+z_T^2]}{2\,R(t)^2}.
  \label{eq:doppler_full}
\end{equation}
For $\theta=0$ ($v_y=0$): $f_D=v_x\Delta x_k/R_0$.

The Doppler is estimated via STFT on the steered beam:
\begin{equation}
  y_k(t) = \frac{1}{N}\sum_{n=0}^{N-1}
  x_{k,n}(t)\;e^{-j\pi n\hat{\ell}_x(t)}.
  \label{eq:steered_beam}
\end{equation}

In the far field, the Doppler rate:
\begin{equation}
  \dot{f}_k \approx \frac{v_x^2}{R_0}
  + \frac{3v_x v_y\,\Delta x_k^{(0)}}{R_0^2} + \cdots
  \label{eq:fdot_k}
\end{equation}

\subsection{Joint Parameter Estimation}
\label{sec:joint_zone_b}

\subsubsection{Step 1: $v_y/R_0^2$ from DOA-rate variation}
From~\eqref{eq:beta_linear_in_Xk}:
$\hat{b}=-v_y/R_0^2$, $\hat{a}=v_x/R_0+v_y x_0/R_0^2$.

\subsubsection{Step 2: Combining with Doppler}
For $\theta\approx 0$ ($\hat{b}\approx 0$):
$|v_x|=|\hat{\dot{f}}|/|\hat{\beta}|$,\;
$R_0=|v_x|/|\hat{\beta}|$.

For $\theta\ne 0$: nonlinear least-squares fit of $4K$
observables $(\hat{\alpha}_k,\hat{\beta}_k,
\hat{f}_0^{(k)},\hat{\dot{f}}_k)$ to the models.  Closed-form
initialisation:
\begin{equation}
  \hat{R}_0 \approx
  \frac{|\hat{\dot{f}}_{\mathrm{avg}}|}
  {\hat{\beta}_{\mathrm{avg}}^2},\qquad
  \hat{v}_y = -\hat{b}\,\hat{R}_0^2,\qquad
  \hat{v}_x = (\hat{a}-\hat{b}\,\hat{x}_0)\hat{R}_0,\qquad
  \hat{x}_0 = \hat{\alpha}_{\mathrm{avg}}\hat{R}_0
  +\bar{X}_{\mathrm{avg}}.
  \label{eq:joint_init}
\end{equation}

\subsubsection{Step 3: $z_T$ initial estimate}
$\hat{z}_T^{(0)}=0$; refined in Zone~A
(Section~\ref{sec:alg3}).

%======================================================================
\section{Zone~A: Gradient Descent with TV Denoising}
\label{sec:zone_a}
%======================================================================

\subsection{Problem Formulation}

Given measured FSSR samples $b_{k,\ell}=|\tilde{y}_k(t_\ell)|^2$,
retrieve the shadow profile:
\begin{equation}
  \hat{\mathbf{P}} = \argmin_{\mathbf{P}}
  \bigl\{\mathcal{F}(\mathbf{P}) + \tau\,\mathcal{R}(\mathbf{P})\bigr\},
  \label{eq:zone_a_problem}
\end{equation}
where
\begin{equation}
  \mathcal{F}(\mathbf{P})
  = \sum_{k,\ell}
  \bigl[b_{k,\ell} - \varepsilon_k(t_\ell;\mathbf{P})\bigr]^2,
  \label{eq:fidelity}
\end{equation}
\begin{equation}
  \mathcal{R}(\mathbf{P}) = \TV(\mathbf{P})
  = \sum_{p,q}\sqrt{(P_{p,q}{-}P_{p+1,q})^2
  +(P_{p,q}{-}P_{p,q+1})^2}.
  \label{eq:tv_reg}
\end{equation}

\subsection{Gradient of the Fidelity Function}
\label{sec:gradient}

Using the combined coefficient~\eqref{eq:H_matrix}:
\begin{equation}
  \frac{\partial\varepsilon_k}{\partial P_{p,q}}
  = 2\,\mathrm{Re}\!\bigl[H_{k,p,q}(t_\ell)\,
  \overline{(1{+}S_{k,\ell})}\bigr].
  \label{eq:deps_dP}
\end{equation}

The full gradient from sample $(k,\ell)$:
\begin{equation}
  \boxed{
  \frac{\partial\mathcal{F}_{k,\ell}}{\partial P_{p,q}}
  = 4[\varepsilon_k{-}b_{k,\ell}]\,
  \mathrm{Re}\!\bigl[H_{k,p,q}(t_\ell)\,
  \overline{(1{+}S_{k,\ell})}\bigr].}
  \label{eq:grad_fidelity}
\end{equation}

In matrix form, the gradient at sample $(k,\ell)$ is:
\begin{equation}
  \nabla_{\mathbf{P}}\mathcal{F}_{k,\ell}
  = 4[\varepsilon_k{-}b_{k,\ell}]\,
  \mathrm{Re}\!\bigl[\overline{(1{+}S_{k,\ell})}\;
  \mathbf{H}_k(t_\ell)\bigr]
  \;\in\;\R^{M_z\times M_x},
  \label{eq:grad_matrix}
\end{equation}
where $[\mathbf{H}_k(t_\ell)]_{p,q}=H_{k,p,q}(t_\ell)$ is the full
$M_z\times M_x$ coefficient matrix.

\paragraph{Efficient computation.}
Although $\mathbf{H}_k$ is $M_z\times M_x$ (not a rank-1 outer
product), each column~$q$ is the product of $F_q^{(k)}$ with a
$G$-vector computed at $R_{\mathrm{eff},q}$.  Column-wise
vectorisation keeps the cost at
$O(M_z M_x)$ per $(k,\ell)$ sample.

\subsection{Algorithm 1: Shadow Profile Retrieval}
\label{sec:alg1}

\begin{algorithm}[H]
\caption{Gradient descent with TV denoising for shadow profile
retrieval}
\label{alg:profile_retrieval}
\begin{algorithmic}[1]
\Require FSSR measurements $\{b_{k,\ell}\}$; coefficient matrices
$\{\mathbf{H}_k(t_\ell)\}$; step size $\gamma>0$;
TV parameter $\tau>0$; denoising iterations~$N_D$
\State Initialise: $\hat{\mathbf{P}}_0\leftarrow\mathbf{0}$,
$\mathbf{s}_0\leftarrow\hat{\mathbf{P}}_0$, $q_0\leftarrow 1$
\For{$t=1,\ldots,N_S$}
  \State $\gamma_t \leftarrow \gamma/\sqrt{t}$
  \State $\mathbf{z}_t \leftarrow \mathbf{s}_{t-1}
    - \frac{\gamma_t}{L_{\mathrm{tot}}}
    \sum_{k,\ell}\nabla_{\mathbf{P}}\mathcal{F}_{k,\ell}
    (\mathbf{s}_{t-1})$
  \Comment{$L_{\mathrm{tot}}=\sum_k L_k$}
  \State $\hat{\mathbf{P}}_t \leftarrow
    \mathrm{prox}_{\mathcal{R}}(\mathbf{z}_t,\gamma_t\tau)$
  \Comment{Algorithm~\ref{alg:tv_denoise}}
  \State $q_t \leftarrow (1{+}\sqrt{1{+}4q_{t-1}^2})/2$
  \State $\mathbf{s}_t \leftarrow \hat{\mathbf{P}}_t
    + \frac{q_{t-1}-1}{q_t}
    (\hat{\mathbf{P}}_t{-}\hat{\mathbf{P}}_{t-1})$
\EndFor
\State \Return $\hat{\mathbf{P}}_{N_S}$
\end{algorithmic}
\end{algorithm}

\subsection{Algorithm 2: TV Denoising Proximal Operator}
\label{sec:alg2}

The proximal operator solves
\begin{equation}
  \mathrm{prox}_{\mathcal{R}}(\mathbf{z},\delta)
  = \argmin_{\mathbf{P}}\Bigl\{
  \tfrac{1}{2\delta}\|\mathbf{P}{-}\mathbf{z}\|^2
  + \TV(\mathbf{P})\Bigr\}
  \quad\text{s.t.}\quad 0\le P_{p,q}\le 1.
  \label{eq:prox_problem}
\end{equation}

\paragraph{Operators.}
\begin{itemize}
  \item \textbf{Discrete gradient}
    $\mathcal{L}^T$:
    $p_{i,j}=m_{i,j}{-}m_{i+1,j}$,\;
    $q_{i,j}=m_{i,j}{-}m_{i,j+1}$.
  \item \textbf{Discrete divergence}
    $\mathcal{L}$:
    $[\mathcal{L}(\mathbf{p},\mathbf{q})]_{i,j}
    = p_{i,j}{+}q_{i,j}{-}p_{i-1,j}{-}q_{i,j-1}$
    (boundary: $p_{0,j}{=}q_{i,0}{=}0$).
  \item \textbf{Clipping}
    $\mathcal{H}_C$:
    $[\mathcal{H}_C(\mathbf{m})]_{i,j}=\min(1,\max(0,m_{i,j}))$.
  \item \textbf{Ball projection}
    $\mathcal{H}_B$:
    $r_{i,j}=p_{i,j}/\max(1,\sqrt{p_{i,j}^2{+}q_{i,j}^2})$,\;
    $s_{i,j}=q_{i,j}/\max(1,\sqrt{p_{i,j}^2{+}q_{i,j}^2})$.
\end{itemize}

\begin{algorithm}[H]
\caption{TV denoising:
$\mathrm{prox}_{\mathcal{R}}(\mathbf{z},\delta)$}
\label{alg:tv_denoise}
\begin{algorithmic}[1]
\Require Pixel image $\mathbf{z}\in\R^{M_z\times M_x}$,
regularisation $\delta$
\State $(\mathbf{p}_0,\mathbf{q}_0)\leftarrow(\mathbf{0},\mathbf{0})$;
$(\mathbf{r}_1,\mathbf{s}_1)\leftarrow(\mathbf{0},\mathbf{0})$;
$w_1\leftarrow 1$
\For{$t=1,\ldots,N_D$}
  \State $(\mathbf{p}_t,\mathbf{q}_t) \leftarrow
    \mathcal{H}_B\!\bigl[(\mathbf{r}_t,\mathbf{s}_t)
    + \tfrac{1}{8\delta}\,\mathcal{L}^T\!\bigl(
    \mathcal{H}_C[\mathbf{z}{-}\delta\,
    \mathcal{L}(\mathbf{r}_t,\mathbf{s}_t)]\bigr)\bigr]$
  \State $w_{t+1}\leftarrow(1{+}\sqrt{1{+}4w_t^2})/2$
  \State $(\mathbf{r}_{t+1},\mathbf{s}_{t+1})\leftarrow
    (\mathbf{p}_t,\mathbf{q}_t)
    +\tfrac{w_t-1}{w_{t+1}}
    [(\mathbf{p}_t{-}\mathbf{p}_{t-1},\,
    \mathbf{q}_t{-}\mathbf{q}_{t-1})]$
\EndFor
\State $\hat{\mathbf{x}}\leftarrow
  \mathcal{H}_C[\mathbf{z}{-}\delta\,
  \mathcal{L}(\mathbf{r}_{N_D},\mathbf{s}_{N_D})]$
\State \Return $\hat{\mathbf{x}}$
\end{algorithmic}
\end{algorithm}

\subsection{Algorithm 3: Joint Motion Parameter and Profile Estimation}
\label{sec:alg3}

The outer loop optimises $\bm{\theta}=(v_x,v_y,z_T,R_0)$ around the
profile retrieval.  The cost function:
\begin{equation}
  C(\hat{\mathbf{P}},\bm{\theta})
  = \sum_{k,\ell}\bigl[b_{k,\ell}
  - \varepsilon_k(t_\ell;\hat{\mathbf{P}},\bm{\theta})\bigr]^2.
  \label{eq:cost_joint}
\end{equation}

Note: changing $\bm{\theta}$ affects the coefficient matrices
$\mathbf{H}_k$ through $R(t_\ell)$, $\Delta x_k(t_\ell)$,
$R_{\mathrm{eff},q}(t_\ell)$ (including the heading angle~$\theta$ in
$\tan\theta$), $F_q^{(k)}$, and $G_{p,q}$.

Numerical gradients for each parameter~$\theta_i$:
\begin{equation}
  \frac{\partial C}{\partial\theta_i}
  \approx
  \frac{C(\hat{\mathbf{P}},\bm{\theta}{+}\Delta\theta_i\mathbf{e}_i)
  - C(\hat{\mathbf{P}},\bm{\theta})}{\Delta\theta_i}.
  \label{eq:num_grad}
\end{equation}

\begin{algorithm}[H]
\caption{Joint shadow profile and motion parameter estimation}
\label{alg:joint_estimation}
\begin{algorithmic}[1]
\Require FSSR measurements $\{b_{k,\ell}\}$; initial estimates
$\hat{v}_x^{(0)},\hat{v}_y^{(0)},\hat{z}_T^{(0)},\hat{R}_0^{(0)},
\hat{x}_0$ (fixed)
\State $q_0\leftarrow 1$;
$\bm{\theta}^{(0)}\leftarrow
(\hat{v}_x^{(0)},\hat{v}_y^{(0)},\hat{z}_T^{(0)},\hat{R}_0^{(0)})$
\For{$n=1,\ldots,N_V$}
  \State Compute trajectories and
    $R_{\mathrm{eff},q}(t_\ell)=R(t_\ell)-X'_q\tan\theta^{(n-1)}$
    \Comment{$\theta=\arctan(v_y/v_x)$}
  \State Build coefficient matrices
    $\mathbf{H}_k(t_\ell)$ using~\eqref{eq:H_matrix}
  \State $\hat{\mathbf{P}}_n \leftarrow
    \text{Algorithm~\ref{alg:profile_retrieval}}
    (\{b_{k,\ell}\},\bm{\theta}^{(n-1)})$
  \State Compute $C$ and numerical gradients~\eqref{eq:num_grad}
  \State $\gamma_i^{(n)}\leftarrow\gamma_i/n$;\quad
    $\tilde{\theta}_i \leftarrow \theta_i^{(n-1)}
    - \gamma_i^{(n)}\,\partial C/\partial\theta_i$
  \State Nesterov:
    $q_n\leftarrow(1{+}\sqrt{1{+}4q_{n-1}^2})/2$;\quad
    $\theta_i^{(n)}\leftarrow\tilde{\theta}_i
    +\tfrac{q_{n-1}-1}{q_n}(\tilde{\theta}_i{-}\theta_i^{(n-1)})$
  \State Enforce $R_0^{(n)}\ge R_{\min}$
\EndFor
\State Final retrieval with converged~$\bm{\theta}$; hard-threshold at
$0.5$
\State \Return
  $\hat{\mathbf{P}},\hat{v}_x,\hat{v}_y,\hat{z}_T,\hat{R}_0$
\end{algorithmic}
\end{algorithm}

%======================================================================
\section{Scale Ambiguity Analysis}
\label{sec:ambiguity}
%======================================================================

\subsection{Single-Receiver Ambiguity}
Under the scaling
$(v_x,v_y,z_T,x_0)\to(\mu v_x,\mu v_y,\mu z_T,\mu x_0)$,
$R_0\to\mu^2 R_0$,
$F'(x',z')=F(x'/\mu,z'/\mu)$,
the FSSR is unchanged~\cite{IETDraft4}.

\subsection{Ambiguity Breaking with Multiple Nodes}
With receivers at known, unscaled positions~$\bar{X}_k$:
$\Delta x_k\to\mu x_T-\bar{X}_k\ne\mu\Delta x_k$ for
$\bar{X}_k\ne 0$.  The ambiguity is broken in Zone~A where
$x_T\sim\bar{X}_k$.

%======================================================================
\section{DOA Estimation Algorithms}
\label{sec:doa_algs}
%======================================================================

\subsection{MUSIC}
Applied to the phase-compensated covariance~\eqref{eq:cov_compensated}.
Noise subspace~$\mathbf{E}_n$; pseudo-spectrum:
\begin{equation}
  P_{\mathrm{MUSIC}}(u)
  = \frac{1}{\mathbf{a}^H(u)\,\mathbf{E}_n\mathbf{E}_n^H\,
  \mathbf{a}(u)},\qquad
  \mathbf{a}(u)=[1,e^{j\pi u},\ldots,e^{j\pi(N-1)u}]^T.
  \label{eq:music}
\end{equation}

\subsection{LS-ESPRIT}
Signal subspace $\mathbf{E}_s$; subarrays
$\mathbf{E}_1=\mathbf{E}_s(1{:}N{-}1,:)$,
$\mathbf{E}_2=\mathbf{E}_s(2{:}N,:)$;
$\bm{\Phi}=\mathbf{E}_1^\dagger\mathbf{E}_2$;
eigenvalue $\lambda_i=e^{j\pi u_i}\Rightarrow u_i=\angle\lambda_i/\pi$.

\subsection{Exact DOA-to-Position Inversion}
\begin{equation}
  \Delta x_k = \frac{R(t)\,\hat{\ell}_x}
  {\sqrt{1-\hat{\ell}_x^2}}.
  \label{eq:exact_inversion}
\end{equation}

%======================================================================
\section{Simulation Parameters}
\label{sec:params}
%======================================================================

\begin{table}[H]
\centering
\caption{Default simulation parameters}
\label{tab:params}
\begin{tabular}{@{}lcl@{}}
\toprule
\textbf{Symbol} & \textbf{Default} & \textbf{Description} \\
\midrule
$K$ & 4 & Receiver nodes \\
$N$ & 16 & Elements per node \\
$R_0$ & $10^4$ & Nominal propagation distance ($\lambda$) \\
$x_0$ & 2000 & Initial target $x$-coordinate \\
$v$ & 1.0 & Target speed ($\lambda$/time) \\
$\theta$ & 0 & Heading angle from $x$-axis \\
$z_T$ & 10 & Constant $z$-offset \\
$M_z{\times}M_x$ & $20{\times}20$ & Pixel grid \\
$2\delta$ & 10 & Pixel side ($\lambda$) \\
$\xi_{\mathrm{th}}$ & 2.0 & Zone threshold \\
SNR & 20\,dB & Signal-to-noise ratio \\
\midrule
$\gamma$ & 5000 & Profile step size \\
$\tau$ & $10^{-6}$ & TV regularisation \\
$N_S$ & 200 & Profile iterations \\
$N_D$ & 20 & Denoising iterations \\
$N_V$ & 20 & Outer parameter iterations \\
\bottomrule
\end{tabular}
\end{table}

%======================================================================
\section{Processing Pipeline Summary}
\label{sec:pipeline}
%======================================================================

\begin{enumerate}
  \item \textbf{FSSR Generation.}
    \begin{enumerate}
      \item For each time step~$t_m$: compute $R(t_m)=R_0-v_y t_m$
        and $\Delta x_k(t_m)=x_0+v_x t_m-\bar{X}_k$.
      \item For each pixel column~$q$: compute
        $R_{\mathrm{eff},q}(t_m)=R(t_m)-X'_q\tan\theta$.
      \item Compute $F_q^{(k)}(t_m)$, $G_{p,q}(t_m)$, and
        $H_{k,p,q}(t_m)$.
      \item True FSSR: $\varepsilon_k(t_m)=|1+\sum_{p,q}
        H_{k,p,q}P_{p,q}|^2$.
      \item Generate noisy signals~\eqref{eq:rx_signal}; measure
        $b_k(t_m)=|\tilde{y}_k(t_m)|^2$~\eqref{eq:fssr_measured}.
    \end{enumerate}

  \item \textbf{Zone Classification.}
    $\xi_k(t)=|\Delta x_k(t)|/\sqrt{R(t)}$; Zone~B/A per node.

  \item \textbf{Zone~B.}
    \begin{enumerate}
      \item Pass~1: uncompensated covariance ($M\le
        M_{\max}^{(\mathrm{unc})}$); rough DOA time series.
      \item Pass~2+: phase-compensated covariance ($M\le
        M_{\mathrm{opt}}$); refined DOA.
      \item Fit $\hat{\beta}_k$ vs.\ $\bar{X}_k$ to extract
        $\hat{a},\hat{b}$.
      \item Doppler via STFT; extract $\hat{\dot{f}}_k$.
      \item Estimate $\hat{v}_x,\hat{v}_y,\hat{R}_0,\hat{x}_0$.
    \end{enumerate}

  \item \textbf{Zone~A.}
    \begin{enumerate}
      \item Collect $b_{k,\ell}=|\tilde{y}_k(t_\ell)|^2$ from each
        node's Zone~A window.
      \item Build coefficient matrices $\mathbf{H}_k(t_\ell)$ using
        $R_{\mathrm{eff},q}=\hat{R}(t_\ell)-X'_q\tan\hat{\theta}$.
      \item Algorithm~\ref{alg:joint_estimation}: joint retrieval of
        $\hat{\mathbf{P}}$ and refined parameters.
    \end{enumerate}

  \item \textbf{Output.}
    $\hat{\mathbf{P}}\in\{0,1\}^{M_z\times M_x}$;\;
    $\hat{v}_x,\hat{v}_y,\hat{z}_T,\hat{R}_0$.
\end{enumerate}

\begin{thebibliography}{9}

\bibitem{Paper1}
X.~Shen and D.~Huang,
``Pixel based shadow profile retrieval in forward scatter radar using
forward scatter shadow ratio,''
\emph{IEEE Access}, vol.~12, pp.~60467--60474, 2024.

\bibitem{Draft5}
X.~Shen and D.~Huang,
``A gradient-based method with denoising for target shadow profile
retrieval in forward scatter radar,''
submitted, 2025.

\bibitem{IETDraft4}
X.~Shen and D.~Huang,
``Retrieving target motion parameters together with target shadow
profile in forward scatter radar,''
submitted to \emph{IET Radar, Sonar \& Navigation}, 2025.

\end{thebibliography}

\end{document}